\documentclass[10pt,a4paper]{amsart}
\usepackage[margin=19mm]{geometry}
\usepackage{amsmath,amssymb,mathtools}
\usepackage[scr=rsfs,cal=euler]{mathalfa}
\usepackage{xfrac}
\usepackage[unicode,hidelinks,bookmarks=false]{hyperref}
\newcommand{\ie}{i.e.\ }
\newcommand{\cf}{cf.\ }
\newcommand{\resp}{respectively\ }
\newcommand{\Subsection}{\subsection}
\providecommand{\nobreakdash}{\nobreak\hbox{-}}
\setlength{\emergencystretch}{2em}
\multlinegap0pt
% ---------------------------------------------------------------------------
% Editorial AMS wrapper prelude. The theorem environments and the mathematical macros below are
% transcribed from the paper's own preamble (cached-originals/source0.tex lines 1--138); the AMS amsart
% class, the named format packages of the pinned TeX Live runtime and this editorial formatting are not
% part of the author's text. No mathematics is changed.
% ---------------------------------------------------------------------------
\usepackage{cleveref}
\newtheorem{definition}{Definition}
\newtheorem{theorem}{Theorem}
\newtheorem{lemma}{Lemma}
\newtheorem{corollary}{Corollary}
\newtheorem{example}{Example}
\newtheorem{proposition}{Proposition}
\newtheorem{remark}{Remark}
\newtheorem{conjecture}{Conjecture}
\newcommand{\R}{\mathbb{R}}
\newcommand{\Prob}{\mathbb{P}}
\newcommand{\EX}{\mathbb{E}}
\newcommand{\Bv}{\mathbf{B}^{[v]}}
\newcommand{\Zv}{\mathbf{Z}^{[v]}}
\newcommand{\Nv}{\mathbf{N}^{[v]}}
\newcommand{\bigzero}{\mbox{\normalfont\Large\bfseries 0}}
\newcommand{\rvline}{\hspace*{-\arraycolsep}\vline\hspace*{-\arraycolsep}}
\renewcommand{\v}[1]{\boldsymbol{#1}}
\newcommand{\Y}{{\mathcal{Y}}}
\newcommand{\A}{{\mathcal{A}}}
\newcommand{\X}{{\mathcal{X}}}
\newcommand{\ind}[1]{\mathbbm{1}\{#1\}}
\newcommand{\ex}[1]{\mathbb{E}\left[#1\right]}
\newcommand{\var}[1]{\mathrm{Var}\left[#1\right]}
\newcommand{\zero}{\textbf{0}}
\newcommand{\B}{\v{B}}
\newcommand{\Z}{\v{Z}}
\newcommand{\N}{\v{N}}
\newcommand{\one}{\v{1}}
\newcommand{\bigbracket}[1]{\bigg(#1\bigg)}
\newcommand{\cs}{\mathcal{S}}
\newcommand{\eA}{\EuScript{A}}
\newcommand{\eR}{\EuScript{R}}
\newcommand{\eY}{\EuScript{Y}}
\newenvironment{rcrevision}{}{}
\newenvironment{aarevision}{}{}
\newenvironment{subrevision}{}{}
\newcommand{\aachg}[1]{#1}
\newcommand{\subchg}[1]{#1}

\begin{document}
\title{Multi-Absorbing Phase-Type Distributions for Right-Censored Competing Risks Data}
\author{Zhihao Qiao, Budhi Surya, Azam Asanjarani, Yoni Nazarathy}
\date{Source: arXiv:2609.19921v1; distributed under CC BY 4.0}
\maketitle
\noindent\emph{Source, version and licence.} \emph{Multi-Absorbing Phase-Type Distributions for Right-Censored Competing Risks Data}, by Zhihao Qiao, Budhi Surya, Azam Asanjarani, Yoni Nazarathy. The source of this editable unit is the e-print of \textbf{version v1} of arXiv:2609.19921, https://arxiv.org/abs/2609.19921v1, whose material is distributed under the Creative Commons Attribution 4.0 International licence, http://creativecommons.org/licenses/by/4.0/, as recorded in the arXiv licence metadata for this paper and shown on that abstract page. The whole of the body below is version v1, the newest version of the paper. Full rights notice: licenses/arxiv-2609.19921-CC-BY-4.0.txt.\par\smallskip
\noindent\textit{Editorial scope.} Counted unit: original Proposition 8 of the article, the statement carried by the source environment \texttt{proposition} with source label \texttt{prop:init-omega}, cached-originals/source0.tex lines 1570--1586, printed as \textbf{Proposition 8} exactly as in the article: the three requirement conditions hold simultaneously for every component k = 1,...,n as soon as the weight parameter omega exceeds the maximum of the explicit piecewise thresholds omega*\_k. It is delivered together with the authors' own complete proof as the single primary proof body, source0.tex lines 1587--1610 (24 lines): the \texttt{proof} environment that belongs to the statement and establishes it in full. No proof marker is added and no mathematics is written, repaired or re-derived; the authors' notation and the printed slips of the source are kept literal. Necessary complete context is retained uncounted, in source order: lines 1537--1540, the source \texttt{equation} with source label \texttt{eq:mu\_k-correction-omega} of the article that the retained passages refer to; lines 1544--1547, the source \texttt{equation} with source label \texttt{eq:scv\_calculated} of the article that the retained passages refer to; lines 1570--1610, the source \texttt{proposition} with source label \texttt{prop:init-omega} of the article that the retained passages refer to. The article's own numbering is reproduced by explicit counter settings: the theorem-like counter of the article is set before each retained passage, so the counted statement prints with the number the article gives it. No \texttt{\textbackslash cite} command occurs in any retained passage: the closure of the retained passages over references and citations was computed, it contains no citation at all, so this unit delivers no bibliography entry and the article's own bibliography data is neither spliced into the bundled source nor redistributed. The source's own class is the class \texttt{article}; the e-print ships \texttt{arxiv.sty}, which is not shipped, loaded or redistributed here, and the wrapper is the AMS \texttt{amsart} class with the article's theorem declarations restated and its macros transcribed from its preamble. The retained passages contain no figure environment and no graphics-inclusion command, so no artwork is redistributed and none of the 1 artwork inclusions shipped by the e-print is used. Every declared body of this unit is an exact, unmodified span of source0.tex: no body is a bounded passage comparison, \texttt{omit} was not used anywhere, and consequently no fidelity receipt is asserted in delivery-evidence.json. The complete concatenated original source is bundled as sources/210816/source0.tex. Omitted from this editable unit and therefore absent from it are source0.tex lines 1--1536; 1541--1543; 1548--1569; 1611--2302. The AMS \texttt{amsart} wrapper, the named format packages of the pinned TeX Live runtime and this editorial source, licence and scope paragraph are editorial formatting only.\par\smallskip
\setcounter{section}{7}
\setcounter{equation}{4}
\setcounter{conjecture}{0}
\setcounter{corollary}{1}
\setcounter{definition}{0}
\setcounter{example}{1}
\setcounter{lemma}{3}
\setcounter{proposition}{7}
\setcounter{remark}{0}
\setcounter{theorem}{0}
\begin{equation}
\label{eq:mu_k-correction-omega}
    \mu_{k} = \hat{\mu}_{k} - \frac{1}{\omega},
\end{equation}

\setcounter{section}{7}
\setcounter{equation}{5}
\setcounter{conjecture}{0}
\setcounter{corollary}{1}
\setcounter{definition}{0}
\setcounter{example}{1}
\setcounter{lemma}{3}
\setcounter{proposition}{7}
\setcounter{remark}{0}
\setcounter{theorem}{0}
\begin{equation}
\label{eq:scv_calculated}
    c_{k}^2 = \frac{\hat{c}_{k}^2(\omega \mu_{k}+1)^2-1}{(\omega \mu_{k})^2}.
\end{equation}

\setcounter{section}{7}
\setcounter{equation}{6}
\setcounter{conjecture}{0}
\setcounter{corollary}{1}
\setcounter{definition}{0}
\setcounter{example}{1}
\setcounter{lemma}{3}
\setcounter{proposition}{7}
\setcounter{remark}{0}
\setcounter{theorem}{0}
\begin{proposition}
\label{prop:init-omega}
Requirements (1)--(3) hold for every $k = 1,\ldots,n$ whenever $\omega > \max_{k} \omega^*_k$, where
\begin{equation}
\label{eq:omega-threshold}
\omega^*_k
=
\frac{1}{\hat{\mu}_k}
\times
\begin{cases}
1, & \hat{c}_k^2 \ge 1,\\[1ex]
\dfrac{1}{\sqrt{\hat{c}_k^2}}, & 0 < \hat{c}_k^2 < \tfrac12,\\[2ex]
\dfrac{1+\sqrt{2\hat{c}_k^2-1}}{1-\hat{c}_k^2}, & \tfrac12 \le \hat{c}_k^2 < 1.
\end{cases}
\end{equation}
Moreover, in the regime $\tfrac12 \le \hat{c}_k^2 < 1$ the threshold is sharp: there are values $\omega$ with $\omega\hat{\mu}_k > 1$ but $\omega < \omega^*_k$ for which $c_k^2 > 1$, violating (3).
\end{proposition}

\setcounter{section}{7}
\setcounter{equation}{7}
\setcounter{conjecture}{0}
\setcounter{corollary}{1}
\setcounter{definition}{0}
\setcounter{example}{1}
\setcounter{lemma}{3}
\setcounter{proposition}{8}
\setcounter{remark}{0}
\setcounter{theorem}{0}
\begin{proof}
Fix $k$ and write $\psi = \omega\hat{\mu}_k$. By \eqref{eq:mu_k-correction-omega}, $\omega\mu_k = \psi - 1$, so requirement (1) is $\psi > 1$, and \eqref{eq:scv_calculated} becomes
\begin{equation}
\label{eq:scv-psi}
c_k^2 = \frac{\hat{c}_k^2\,\psi^2 - 1}{(\psi-1)^2}.
\end{equation}
Hence requirement (2) reads $\psi > 1/\sqrt{\hat{c}_k^2}$, and
\[
c_k^2 - 1 = \frac{g(\psi)}{(\psi-1)^2},
\qquad
g(\psi) := (\hat{c}_k^2-1)\,\psi^2 + 2\psi - 2 .
\]
We consider the three regimes of \eqref{eq:omega-threshold}.

If $\hat{c}_k^2 \ge 1$ then for $\psi > 1$ we have $g(\psi) \ge 2\psi - 2 > 0$, so $c_k^2 > 1 > 0$: requirements (2) and (3) are automatic and only (1) binds, giving $\omega^*_k = 1/\hat{\mu}_k$.

If $\hat{c}_k^2 < 1$ then $g$ is a concave parabola with $g(1) = \hat{c}_k^2 - 1 < 0$ and discriminant $4(2\hat{c}_k^2-1)$. When $\hat{c}_k^2 < \tfrac12$ the discriminant is negative, so $g < 0$ everywhere and $c_k^2 < 1$ for every $\psi$: requirement (3) is automatic, and the binding requirement is (2), $\psi > 1/\sqrt{\hat{c}_k^2}$, which also implies (1) since $1/\sqrt{\hat{c}_k^2} > 1$.

When $\tfrac12 \le \hat{c}_k^2 < 1$ the parabola $g$ has real roots
\[
\psi_\pm = \frac{1 \pm \sqrt{2\hat{c}_k^2-1}}{1-\hat{c}_k^2},
\]
with $g > 0$ on $(\psi_-,\psi_+)$ and $g < 0$ outside $[\psi_-,\psi_+]$. For $\psi > \psi_+$ we therefore have $c_k^2 < 1$, which is (3). For (2), note that at $\psi = \psi_+$ equation \eqref{eq:scv-psi} gives $c_k^2 = 1$, so its numerator satisfies $\hat{c}_k^2\psi_+^2 - 1 = (\psi_+-1)^2 > 0$; since the numerator is increasing in $\psi$, it remains positive for all $\psi \ge \psi_+$, giving $c_k^2 > 0$. For (1), $\psi_+ \ge 1/(1-\hat{c}_k^2) > 1$. Hence $\omega^*_k = \psi_+/\hat{\mu}_k$ suffices. For sharpness, $g(1) < 0$ places $1$ outside the interval $(\psi_-,\psi_+)$, so $1 < \psi_- < \psi_+$; any $\psi \in (\psi_-,\psi_+)$ satisfies (1) yet has $g(\psi) > 0$, i.e.\ $c_k^2 > 1$, violating (3).
\end{proof}

\end{document}
